 Classical mechanics and statistical mechanics admit a unified formulation in
 which the fundamental object is not an individual trajectory but a state of
 knowledge about the physical system. From this perspective, statistical
 mechanics is not necessarily a different dynamical theory from mechanics. It
 is an extension of mechanics from individual states to probability
 distributions over individual states.

There is, however, an important distinction between two questions:
\begin{enumerate}
   \item What is the physical state of a system?
   \item What is the mathematical description of our knowledge of that state?
\end{enumerate}
Bayesian reasoning makes this distinction particularly visible.

\section{ Mechanics as a theory of pure states}
Consider an ordinary Hamiltonian system with \(n\) degrees of freedom.
Its phase space is
\[\Gamma=T^*Q,\]
with coordinates
\[z=(q^1,\ldots,q^n,p_1,\ldots,p_n).\]
In classical mechanics, the state at time \(t\) is a point
\[z(t)\in\Gamma.\]
Its evolution is determined by Hamilton's equations:
\[
\dot q^i=\frac{\partial H}{\partial p_i},
\qquad
\dot p_i=-\frac{\partial H}{\partial q^i}.
\]
Equivalently,
\[\dot z=X_H(z).\]
This is what we might call a pure mechanical state: the system is assigned a
definite point in phase space.

Now introduce a probability density on phase space:
\[\rho(q,p,t)\geq 0, \qquad \int_\Gamma \rho\,d\Gamma=1.\]
This describes a statistical state. Instead of specifying one phase-space
point, we specify a probability distribution over possible phase-space points.
The evolution equation is the Liouville equation:
\[\boxed{\frac{\partial\rho}{\partial t}+\{\rho,H\}=0.}\]
Here is the important observation:

\emph{The Liouville equation is not a new dynamical law added to Hamiltonian
mechanics. It is Hamiltonian mechanics lifted from individual points to
probability distributions over points.}
Indeed, let \(\Phi_t\) denote the Hamiltonian flow. Then
\[\rho_t=(\Phi_t)_*\rho_0,\]
where the star denotes the pushforward of the probability measure.

For a pure state,
\[\rho(z,t)=\delta(z-z(t)),\]
the Liouville equation reduces to the Hamiltonian equations for \(z(t)\).
Thus, in this precise sense, ordinary mechanics is a special case of classical
statistical mechanics.

\section{ The Bayesian interpretation}
Suppose we do not know the exact initial phase-space point. We have information
about it, represented by a probability distribution
\[\rho_0(z).\]
The distribution is then our state of knowledge concerning the actual
mechanical state.

The distinction between a point and a distribution becomes:
\begin{itemize}
\item Pure state: complete specification of the phase-space point.
\item Mixed state: incomplete specification, represented by a probability
    measure over phase space.
\end{itemize}
This is analogous to the distinction between pure and mixed states in quantum
mechanics, although there are important differences.
For Bayesian inference, probabilities quantify uncertainty or degrees of
belief, conditional on the available information.
The probability distribution is therefore not necessarily a statement that the
system itself is physically spread over phase space. It may simply represent
our ignorance of its actual state.

For example, consider a harmonic oscillator.
Its actual state may be
\[(q_0,p_0),\]
but our knowledge might be described by a Gaussian distribution centered at that point:
\[
\rho(q,p)=
\frac{1}{2\pi\sigma_q\sigma_p}
\exp\left[
-\frac{(q-q_0)^2}{2\sigma_q^2}
-\frac{(p-p_0)^2}{2\sigma_p^2}
\right].
\]
Hamiltonian evolution transports this distribution through phase space.
Notice that the underlying dynamics has not changed. Only the description of
the initial state has changed.

\section{The deeper mathematical structure: a convex space of states}
Let us define the space of all normalized probability measures on phase space:
\[\mathcal S(\Gamma)=\left\{\mu:\mu(\Gamma)=1,\quad \mu\geq 0\right\}.\]
This is a convex space.
If \(\mu_1\) and \(\mu_2\) are two statistical states, then
\[\mu=\lambda\mu_1+(1-\lambda)\mu_2,\qquad 0\leq\lambda\leq1,\]
is also a statistical state.
The crucial mathematical fact is that the extremal elements of this convex
space are precisely the Dirac measures:
\[\mu_z=\delta_z.\]
These are the pure states in the classical statistical description.

Every probability measure can, in fact, be represented as a mixture of these
extremal states:
\[\boxed{\mu=\int_\Gamma \delta_z\,d\mu(z).}\]
This gives us a particularly clean formulation:
\begin{table}
    \begin{tabular}{rr}
        Classical mechanics	& Classical statistical mechanics\\
        Point \(z\in\Gamma\) &	Probability measure \(\mu\) on \(\Gamma\)\\
        Pure state&	Pure or mixed statistical state\\
        Hamiltonian flow&	Liouville evolution\\
        Observable \(A(z(t))\)	&Expectation \(\langle A\rangle_\mu\)=\(\int_\Gamma A(z)\,d\mu_t(z)\)
    \end{tabular}
\end{table}
The evolution of expectations follows directly:
\[\frac{d}{dt}\langle A\rangle=\langle\{A,H\}\rangle.\]
This is the statistical version of Hamiltonian dynamics.

There is a useful parallel here with quantum mechanics. In quantum theory, the
state space is also convex, and pure states are its extremal elements. But the
quantum state space is not simply a space of probability distributions over
classical phase-space points. Its noncommutative structure makes the
distinction fundamental. In classical mechanics, by contrast, the underlying
state space remains the ordinary phase space.

\section{Statistical mechanics is not defined by an infinite number of particles}
Statistical mechanics is often associated with systems containing an enormous
number of particles, but this is not its defining characteristic.
Consider just one particle moving in a potential.
We can formulate statistical mechanics for it:
\begin{itemize}
\item A microcanonical distribution at fixed energy.
\item A canonical distribution at fixed temperature.
\item A general non-equilibrium probability distribution.
\end{itemize}
Conversely, a system with a very large number of particles can be described by
ordinary mechanics if we specify all its initial positions and momenta.
Thus, the distinction is not:
\[
\text{few particles}\quad\longleftrightarrow\quad
\text{many particles}.
\]
Rather, it is:
\[
\boxed{
\text{individual mechanical state}
\quad\longleftrightarrow\quad
\text{probability distribution of mechanical states}.
}
\]
The thermodynamic limit,
\[
N\rightarrow\infty,\qquad
V\rightarrow\infty,\qquad
\frac NV=\text{constant},
\]
is important for explaining phase transitions, thermodynamic singularities and
the emergence of macroscopic behavior.
But it is not required to define statistical mechanics.

\section{Entropy and coarse-graining}
If we consider only the exact microscopic dynamics and probability
distributions over its states, then the difference between mechanics and
statistical mechanics really is largely a difference in the states being
considered.
However, statistical mechanics also introduces concepts that are not contained
in the mere replacement of points by distributions.
The most important is entropy.

For a continuous probability density,
\[S_{\mathrm G}[\rho]=-k_B\int_\Gamma \rho\ln\rho\,d\Gamma,\]
where the reference phase-space measure and the units used inside the logarithm
must be specified.
Under exact Hamiltonian evolution, Liouville's theorem implies
\[\boxed{\frac{dS_{\mathrm G}}{dt}=0.}\]
This is a very significant result.
The microscopic statistical state evolves reversibly, just as an individual
mechanical state does.

Yet thermodynamics tells us that entropy increases.
Where does the difference come from?
It comes from the fact that thermodynamic entropy is not simply the
fine-grained Gibbs entropy.
One introduces coarse-graining, reduced descriptions, macroscopic variables, or
other restrictions on the information retained.
For example, a distribution may develop increasingly fine structures in phase
space. Its fine-grained entropy remains constant, but its coarse-grained
entropy can increase.
This is one of the fundamental connections between reversible mechanics and
irreversible thermodynamics.

It also means that one should distinguish three levels:
\begin{enumerate}
   \item Microscopic mechanics: evolution of individual phase-space points.
   \item Fine-grained statistical mechanics: evolution of probability distributions.
   \item  Macroscopic statistical physics: evolution of reduced descriptions and thermodynamic variables.
\end{enumerate}
The first two are related by a mathematically precise extension of the same
dynamics. The third involves additional conceptual and mathematical
ingredients.

\section{Bayesian updating versus dynamical evolution}
There is another distinction worth emphasizing.
In Bayesian theory, a probability distribution can change for two different reasons.

\paragraph{A. Dynamical evolution}
We know that the physical system evolves, and therefore transport our
probability distribution according to Hamiltonian mechanics:
\[\rho_t=(\Phi_t)_*\rho_0.\]

\paragraph{B. Acquisition of information}
We obtain new information, for example through a measurement, and update the
distribution:
\[\rho(z\mid D)=\frac{P(D\mid z)\rho(z)}{\int P(D\mid z')\rho(z')\,dz'}.\]
These are conceptually different operations.
The first describes the physical evolution of the system's state of knowledge
under a known dynamical law.
The second describes a change in our knowledge resulting from new evidence.

% do not understand that!@!!!!
This distinction is especially important when discussing statistical mechanics
from an epistemic or Bayesian standpoint.
For instance, equilibrium distributions need not mean that the microscopic
system has stopped evolving. They can represent stationary probability
distributions, or equilibrium descriptions obtained after information has been
discarded.

\section{ The difficult question: is a mixed state merely ignorance?}
There are at least three interpretations of a statistical state.

\paragraph{Epistemic interpretation}
The system has a definite microscopic state \(z\), but we do not know it. The
distribution represents our uncertainty.

\paragraph{Ensemble interpretation}
The distribution describes a collection of similarly prepared systems, each
having its own definite phase-space point.

\paragraph{Objective statistical interpretation}
The probability distribution is itself regarded as a physical description,
possibly because the microscopic state is inaccessible or because the relevant
physical description is intrinsically statistical.

These interpretations can give identical mathematical predictions for many
experiments, but they are not conceptually equivalent.

There is also a subtle issue concerning ergodicity.
Suppose the system is deterministic and has a definite initial point. Its
trajectory might explore a large part of an energy surface. Under suitable
ergodic assumptions, long-time averages can equal ensemble averages:
\[\lim_{T\to\infty}\frac1T\int_0^T A(z(t))\,dt=
\int_\Gamma A(z)\rho_{\rm eq}(z)\,d\Gamma.\]
But this does not mean that the individual trajectory has become a probability
distribution.
It means that two different procedures for calculating averages can agree.
This distinction is important: the probability description does not
automatically replace the underlying mechanics.

\section{The hierarchy of descriptions}
For a system of \(N\) particles, the full mechanical phase space is
\[\Gamma_N=T^*\mathbb R^{3N}.\]
A statistical description is a probability density
\[\rho_N(\mathbf q_1,\ldots,\mathbf q_N,\mathbf p_1,\ldots,\mathbf p_N).\]
But in practical statistical mechanics, we rarely work with this complete
distribution.
Instead, we introduce reduced distributions:
\[f_1(\mathbf q,\mathbf p),\]
\[f_2(\mathbf q_1,\mathbf p_1;\mathbf q_2,\mathbf p_2),\]
and so on.
These are obtained by integrating out degrees of freedom.
The evolution of these distributions leads to the BBGKY hierarchy.

This gives us a sequence of descriptions:
\begin{equation*}
    \begin{gathered}
\text{individual microscopic states}\\
\downarrow \\
\text{full probability distribution}\\
\downarrow \\
\text{reduced distributions}\\
\downarrow \\
\text{macroscopic variables}.
    \end{gathered}
\end{equation*}
The important point is that the successive levels are not necessarily different
physical theories. They are different levels of description of the same
underlying dynamics, although approximations introduced in passing between them
can produce new effective dynamical laws.
This is particularly relevant to kinetic theory, where the transition from the
Liouville equation to the Boltzmann equation involves assumptions such as
molecular chaos.

Classical mechanics and classical statistical mechanics are two levels of
description within a single dynamical framework. The fundamental dynamical
structure is Hamiltonian mechanics on phase space. Its pure states are points
of phase space, while its general states are probability measures over phase
space. Hamiltonian dynamics induces a natural evolution on the entire convex
space of statistical states.

In mathematical terms:
\[
\boxed{
\begin{aligned}
&\text{Configuration of the system:} &&\Gamma,\\
&\text{Pure states:} &&z\in\Gamma,\\
&\text{General states:} &&\rho\in\mathcal S(\Gamma),\\
&\text{Dynamics:} &&\Phi_t:\Gamma\to\Gamma,\\
&\text{Statistical dynamics:} &&
(\Phi_t)_*:\mathcal S(\Gamma)\to\mathcal S(\Gamma).
\end{aligned}
}
\]
This is a very natural generalization of the usual Hamiltonian formulation.
There is also a useful connection with the Liouville formulation of mechanics,
in which pure and mixed states can be treated within the same formalism.

\section{Comparison with quantum mechanics}
We can put the three descriptions side by side:
\begin{table}
    \begin{tabular}{ccc}
        &Classical&	Quantum\\
        Pure state&	Phase-space point \(z\)	&State vector \(\psi\) \\
        General state&	Probability measure \(\rho\)&	Density operator \(\hat\rho\)\\
        Observable	&Function \(A(z)\)&	Operator \(\hat A\)\\
        Expectation	&\(\int A\rho\,d\Gamma\)&	\(\mathrm{Tr}(\hat\rho\hat A)\)\\
        Evolution&	Liouville equation&	von Neumann equation\\
    \end{tabular}
\end{table}
The quantum evolution is
\[i\hbar\frac{d\hat\rho}{dt}=[\hat H,\hat\rho].\]
For a pure quantum state,
\[\hat\rho=|\psi\rangle\langle\psi|,\]
this reduces to the Schrödinger equation (up to the irrelevant overall phase).

Notice the structural similarity:
\begin{itemize}
   \item Classical mechanics is the pure-state sector of classical statistical mechanics.
   \item Quantum mechanics is the pure-state sector of quantum statistical mechanics.
\end{itemize}
But there is an essential difference: classical statistical states are
probability measures on phase space, whereas quantum statistical states are
positive operators on Hilbert space. Quantum superposition, noncommutativity
and entanglement have no straightforward classical counterpart.

 The distinction between mechanics and statistical mechanics is not
 fundamentally a distinction between few and infinitely many particles, nor
 necessarily between two different dynamical theories. At the microscopic
 level, it can be understood as a distinction between pure and mixed states of
 one and the same dynamical system.
The genuinely additional content of statistical physics enters when we ask how
to obtain macroscopic descriptions, equilibrium, irreversibility and
thermodynamics from that unified microscopic framework.

